\documentclass[11pt,letterpaper]{article}

% Packages
\usepackage[utf8]{inputenc}
\usepackage[T1]{fontenc}
\usepackage{geometry}
\usepackage{graphicx}
\usepackage{booktabs}
\usepackage{longtable}
\usepackage{array}
\usepackage{multirow}
\usepackage{xcolor}
\usepackage{colortbl}
\usepackage{hyperref}
\usepackage{cleveref}
\usepackage{amsmath}
\usepackage{amssymb}
\usepackage{siunitx}
\usepackage{natbib}
\usepackage{lineno}
\usepackage{setspace}
\usepackage{enumitem}
\usepackage{tikz}
\usepackage{pdflscape}

% Page setup
\geometry{margin=1in}
\onehalfspacing
\linenumbers

% Custom colors
\definecolor{complete}{RGB}{107,203,119}
\definecolor{inprogress}{RGB}{255,217,61}
\definecolor{pending}{RGB}{255,107,107}
\definecolor{high}{RGB}{46,125,50}
\definecolor{moderate}{RGB}{255,152,0}
\definecolor{low}{RGB}{198,40,40}

% Custom commands
\newcommand{\status}[1]{\textcolor{#1}{\textbullet}}
\newcommand{\HIGH}{\textcolor{high}{\textbf{HIGH}}}
\newcommand{\MODERATE}{\textcolor{moderate}{\textbf{MODERATE}}}
\newcommand{\LOW}{\textcolor{low}{\textbf{LOW}}}
\newcommand{\variable}[1]{\texttt{#1}}
\newcommand{\unit}[1]{\ensuremath{\mathrm{#1}}}

% Title
\title{%
\textbf{Master Planning Document}\\[0.5em]
\Large Remote Sensing Detection of Flood-Buried Forest Deposits:\\
A Framework for Quantifying Cryptic Carbon Reservoirs\\
in Amazonian Floodplains\\[1em]
\large CICRA Pilot Study, Madre de Dios, Peru
}

\author{%
[Authors for Expert Revision]\\
\textit{Draft Version 1.0}
}

\date{April 2026}

\begin{document}

\maketitle

\begin{abstract}
This master planning document consolidates the research framework, variable definitions, data product comparisons, and prediction model alternatives for detecting and quantifying flood-buried forest carbon deposits using remote sensing. The project addresses critical knowledge gaps identified by the NASA Earth Science Decadal Survey (2017--2027) and National Academies regarding terrestrial carbon stocks and residence times. We present a multi-platform approach comparing digital elevation models, spaceborne lidar, synthetic aperture radar, and hyperspectral imagery for terrace scarp detection and buried wood identification in the Madre de Dios River basin, Peru.
\end{abstract}

\tableofcontents
\newpage

%==============================================================================
\section{Motivation and Knowledge Gaps}
\label{sec:motivation}
%==============================================================================

\subsection{The Cryptic Carbon Problem}

Global carbon cycle models struggle to balance terrestrial carbon budgets, with significant ``missing sinks'' partially attributable to unquantified burial pathways. The terrestrial carbon cycle remains the \textbf{least constrained component} of the global carbon budget:

\begin{itemize}[noitemsep]
    \item Terrestrial CO$_2$ uptake varies from a net \SI{4.1}{Pg~C~yr^{-1}} sink to a \SI{0.4}{Pg~C~yr^{-1}} source
    \item This variability accounts for the majority of interannual atmospheric CO$_2$ growth fluctuations
    \item Large uncertainties stem from poor understanding of carbon allocation, stocks, and residence times
\end{itemize}

\subsection{NASA Decadal Survey Alignment (2017--2027)}

The project directly addresses priorities from ``Thriving on Our Changing Planet'':

\begin{table}[h]
\centering
\caption{Alignment with NASA Decadal Survey Priority Science Questions}
\label{tab:decadal_alignment}
\begin{tabular}{p{5cm}p{6cm}c}
\toprule
\textbf{Decadal Survey Question} & \textbf{Project Contribution} & \textbf{Priority} \\
\midrule
What are the structure, function, and biodiversity of Earth's ecosystems? & Maps terrace morphology indicating buried forest structure & \HIGH \\
\addlinespace
How is the water cycle changing? & Documents historical forest burial by flooding events & \MODERATE \\
\addlinespace
How do ecosystems respond to human actions? & Quantifies cryptic carbon not in current budgets & \HIGH \\
\bottomrule
\end{tabular}
\end{table}

\subsection{National Academies Research Priorities}

From the \textit{Negative Emissions Technologies and Reliable Sequestration} report (2019):

\begin{enumerate}[noitemsep]
    \item Terrestrial carbon sequestration = increase/maintenance of organic C in biological stocks over \textbf{multidecadal timescales}
    \item Stocks of interest: woody biomass, coarse woody debris, soil organic matter
    \item Global wood burial potential: \SI{10 \pm 5}{Gt~C~yr^{-1}}
    \item Tropical forests have largest potential: \SI{4.2}{Gt~C~yr^{-1}}
    \item Currently $\sim$\SI{65}{Gt~C} on forest floors as coarse woody debris globally
\end{enumerate}

\subsection{Key Knowledge Gaps Addressed}

\begin{table}[h]
\centering
\caption{Knowledge Gaps and Project Contributions}
\label{tab:knowledge_gaps}
\begin{tabular}{p{3.5cm}p{5cm}p{4cm}}
\toprule
\textbf{Knowledge Gap} & \textbf{Description} & \textbf{Project Response} \\
\midrule
Subsurface carbon stocks & Most evidence focused on above-ground biomass and top \SI{1}{m} soil & Focus on meters-deep burial deposits \\
\addlinespace
Carbon residence times & Poor understanding of long-term storage & Wood burial extends residence to centuries--millennia \\
\addlinespace
Floodplain carbon storage & Poorly quantified & Remote sensing framework for mapping \\
\addlinespace
Episodic burial events & Megafloods not in disturbance models & Links scarp morphology to burial events \\
\bottomrule
\end{tabular}
\end{table}

%==============================================================================
\section{Study Area and Sampling Design}
\label{sec:study_area}
%==============================================================================

\subsection{Madre de Dios River, Peru}

\begin{table}[h]
\centering
\caption{Study Area Characteristics}
\label{tab:study_area}
\begin{tabular}{ll}
\toprule
\textbf{Parameter} & \textbf{Value} \\
\midrule
Location & \ang{12.57}S, \ang{70.10}W \\
Elevation range & \SIrange{200}{310}{m} asl \\
Annual precipitation & $\sim$\SI{2500}{mm} \\
Mean river discharge & $\sim$\SI{3000}{m^3~s^{-1}} \\
Geomorphic setting & Andean foreland basin \\
Primary study area & \SI{8 x 6}{km} (CICRA station) \\
\bottomrule
\end{tabular}
\end{table}

\subsection{Sampling Locations}

\begin{table}[h]
\centering
\caption{Ground Truth Sampling Sites}
\label{tab:sampling_sites}
\begin{tabular}{lrrp{4cm}}
\toprule
\textbf{Site} & \textbf{Longitude} & \textbf{Latitude} & \textbf{Description} \\
\midrule
COLUMNA\_1\_CICRA & \ang{-70.1049} & \ang{-12.5653} & CICRA terrace exposure \\
COLUMNA\_4\_CICRA & \ang{-70.1019} & \ang{-12.5690} & CICRA terrace exposure \\
COLUMNA\_5\_LOS\_AMIGOS & \ang{-70.0906} & \ang{-12.5572} & Los Amigos confluence \\
TOBA\_PATRICE\_PM & \ang{-69.2053} & \ang{-12.5711} & Puerto Maldonado area \\
TRONCOS\_PM & \ang{-69.1900} & \ang{-12.5786} & Fossil logs exposure \\
\bottomrule
\end{tabular}
\end{table}

%==============================================================================
\section{Variables of Interest}
\label{sec:variables}
%==============================================================================

\subsection{Primary Response Variables}

\begin{table}[h]
\centering
\caption{Primary Response Variables}
\label{tab:response_vars}
\begin{tabular}{p{3cm}lllp{3.5cm}}
\toprule
\textbf{Variable} & \textbf{Symbol} & \textbf{Units} & \textbf{Range} & \textbf{Source} \\
\midrule
Scarp probability & $P_{scarp}$ & -- & 0--1 & DEM derivatives \\
Wood presence & $P_{wood}$ & -- & 0--1 & Spectral + field \\
Carbon stock & $C_{stock}$ & \unit{Mg~C} & TBD & Integrated model \\
\bottomrule
\end{tabular}
\end{table}

\subsection{Terrain Predictor Variables}

\begin{table}[h]
\centering
\caption{Terrain-Derived Predictor Variables}
\label{tab:terrain_vars}
\begin{tabular}{p{3.5cm}llll}
\toprule
\textbf{Variable} & \textbf{Symbol} & \textbf{Units} & \textbf{Range} & \textbf{Method} \\
\midrule
Slope & $\theta$ & degrees & 0--35 & Horn's algorithm \\
Aspect & $\alpha$ & degrees & 0--360 & 8-neighbor \\
Topographic Position Index & $TPI_r$ & m & $\pm$15 & Multi-scale ($r$ = 3,5,10,15 cells) \\
Profile curvature & $\kappa_p$ & \unit{m^{-1}} & $\pm$0.05 & 2nd derivative \\
Plan curvature & $\kappa_c$ & \unit{m^{-1}} & $\pm$0.05 & 2nd derivative \\
Terrain Ruggedness Index & $TRI$ & m & 0--10 & Riley et al. \\
Edge magnitude & $E_{mag}$ & -- & 0--1 & Sobel filter \\
TPI variance & $\sigma^2_{TPI}$ & \unit{m^2} & 0--50 & Multi-scale \\
\bottomrule
\end{tabular}
\end{table}

\subsection{Spectral Predictor Variables}

\begin{table}[h]
\centering
\caption{Spectral Index Variables}
\label{tab:spectral_vars}
\begin{tabular}{p{4cm}lll}
\toprule
\textbf{Index} & \textbf{Formula} & \textbf{Range} & \textbf{Purpose} \\
\midrule
NDVI & $\frac{NIR - Red}{NIR + Red}$ & $-1$ to $+1$ & Vegetation density \\
\addlinespace
NDWI & $\frac{Green - NIR}{Green + NIR}$ & $-1$ to $+1$ & Moisture content \\
\addlinespace
NBR & $\frac{NIR - SWIR_2}{NIR + SWIR_2}$ & $-1$ to $+1$ & Bare/burned surface \\
\addlinespace
SOCI & $\frac{Red - Blue}{Red + Blue}$ & $-1$ to $+1$ & Soil organic carbon \\
\addlinespace
BSI & $\frac{(SWIR + Red) - (NIR + Blue)}{(SWIR + Red) + (NIR + Blue)}$ & $-1$ to $+1$ & Bare soil \\
\addlinespace
Cellulose Index & Absorption at \SI{2100}{nm} & 0--1 & Wood detection \\
\addlinespace
Lignin Index & Absorption at \SI{2270}{nm} & 0--1 & Organic matter \\
\bottomrule
\end{tabular}
\end{table}

\subsection{Lidar-Derived Variables}

\begin{table}[h]
\centering
\caption{Lidar-Derived Variables}
\label{tab:lidar_vars}
\begin{tabular}{p{3.5cm}llll}
\toprule
\textbf{Variable} & \textbf{Symbol} & \textbf{Units} & \textbf{Accuracy} & \textbf{Source} \\
\midrule
Ground elevation & $z_{ground}$ & m & $\pm$\SI{1}{m} & GEDI L2A \\
Canopy height (RH100) & $h_{canopy}$ & m & $\pm$\SI{4}{m} & GEDI L2A \\
Aboveground biomass density & $AGBD$ & \unit{Mg~ha^{-1}} & $\pm$20\% & GEDI L4A \\
Scarp height & $H_{scarp}$ & m & $\pm$\SI{2}{m} & GEDI + DEM \\
\bottomrule
\end{tabular}
\end{table}

\subsection{Carbon Estimation Parameters}

\begin{table}[h]
\centering
\caption{Carbon Quantification Parameters}
\label{tab:carbon_params}
\begin{tabular}{p{4cm}lllp{3cm}}
\toprule
\textbf{Parameter} & \textbf{Symbol} & \textbf{Units} & \textbf{Value/Range} & \textbf{Source} \\
\midrule
Total scarp length & $L_{total}$ & km & TBD & DEM vectorization \\
Mean scarp height & $H_{mean}$ & m & 5--15 & GEDI + field \\
Volume per meter & $V_{m}$ & \unit{m^3~m^{-1}} & 0.5--2.0 & Field measurement \\
Wood density (subfossil) & $\rho_{wood}$ & \unit{kg~m^{-3}} & 400--700 & Lab analysis \\
Carbon fraction & $f_C$ & -- & 0.47 & Literature \\
Preservation factor & $f_{pres}$ & -- & 0.50--0.97 & Lab analysis \\
\bottomrule
\end{tabular}
\end{table}

\subsubsection{Carbon Stock Equation}

The total carbon stock is estimated as:

\begin{equation}
\label{eq:carbon_stock}
C_{stock} = L_{total} \times P_{wood} \times V_m \times \rho_{wood} \times f_C \times f_{pres}
\end{equation}

With uncertainty propagation:

\begin{equation}
\label{eq:uncertainty}
\frac{\sigma_C^2}{C^2} = \left(\frac{\sigma_L}{L}\right)^2 + \left(\frac{\sigma_P}{P}\right)^2 + \left(\frac{\sigma_V}{V}\right)^2 + \left(\frac{\sigma_\rho}{\rho}\right)^2 + \left(\frac{\sigma_f}{f}\right)^2
\end{equation}

%==============================================================================
\section{Data Products and Platform Comparison}
\label{sec:platforms}
%==============================================================================

\subsection{Digital Elevation Models}

\begin{table}[h]
\centering
\caption{DEM Product Comparison for Scarp Detection}
\label{tab:dem_comparison}
\small
\begin{tabular}{p{2.5cm}cccccc}
\toprule
\textbf{Product} & \textbf{Resolution} & \textbf{Vert. Acc.} & \textbf{Forest Perf.} & \textbf{Access} & \textbf{Scarp Det.} & \textbf{Priority} \\
\midrule
SRTM v3 & \SI{30}{m} & $\pm$\SI{9}{m} & Poor & Free & $\star\star$ & Low \\
NASADEM & \SI{30}{m} & $\pm$\SI{5}{m} & Moderate & Free & $\star\star\star$ & Medium \\
Copernicus GLO-30 & \SI{30}{m} & $\pm$\SI{4}{m} & Moderate & Free & $\star\star\star$ & Medium \\
\rowcolor{complete!30}
FABDEM & \SI{30}{m} & $\pm$\SI{2-3}{m} & \textbf{Excellent} & Free & $\star\star\star\star$ & \textbf{HIGH} \\
TanDEM-X/WorldDEM & \SI{12}{m} & $\pm$\SI{2}{m} & Moderate & Commercial & $\star\star\star\star$ & Low \\
WorldDEM DTM & \SI{12}{m} & $\pm$\SI{2}{m} & Good & Commercial & $\star\star\star\star\star$ & Low \\
\bottomrule
\end{tabular}
\end{table}

\textbf{Key Finding:} FABDEM removes forest canopy bias, roughly halving vertical error in forested regions compared to Copernicus GLO-30.

\subsection{Spaceborne Lidar}

\begin{table}[h]
\centering
\caption{Spaceborne Lidar Comparison}
\label{tab:lidar_comparison}
\begin{tabular}{p{2.5cm}ccccc}
\toprule
\textbf{Product} & \textbf{Footprint} & \textbf{Ground Acc.} & \textbf{Canopy Acc.} & \textbf{Coverage} & \textbf{Status} \\
\midrule
GEDI L2A & \SI{25}{m} & $\pm$\SI{1}{m} & $\pm$\SI{4-5}{m} & $\pm$\ang{51.6} lat & Active \\
GEDI L4A AGB & \SI{25}{m} & N/A & AGB & $\pm$\ang{51.6} lat & Active \\
ICESat-2 ATL08 & \SI{11 x 100}{m} & $\pm$\SI{0.5-0.7}{m} & $\pm$\SI{1.4-4}{m} & Global & Active \\
\bottomrule
\end{tabular}
\end{table}

\textbf{Tropical Forest Note:} GEDI preferred for dense tropical canopy; ICESat-2 struggles with multi-layered canopy (NRMSE 67\%).

\subsection{Synthetic Aperture Radar}

\begin{table}[h]
\centering
\caption{SAR Mission Comparison for Forest Structure}
\label{tab:sar_comparison}
\begin{tabular}{p{2.5cm}ccccc}
\toprule
\textbf{Mission} & \textbf{Band} & \textbf{Resolution} & \textbf{Biomass Sens.} & \textbf{Access} & \textbf{Status} \\
\midrule
Sentinel-1 & C-band & \SI{10-20}{m} & $<$\SI{50}{Mg~ha^{-1}} & Free & Active \\
ALOS-2 PALSAR-2 & L-band & \SI{25}{m} & $<$\SI{100}{Mg~ha^{-1}} & Free & Active \\
\rowcolor{inprogress!30}
NISAR & L+S-band & \SI{12}{m} & $<$\SI{150}{Mg~ha^{-1}} & Free & July 2025 \\
\rowcolor{complete!30}
ESA BIOMASS & P-band & \SI{50-200}{m} & $>$\SI{300}{Mg~ha^{-1}} & Free & \textbf{Launched} \\
\bottomrule
\end{tabular}
\end{table}

\textbf{Key Finding:} ESA BIOMASS P-band (\SI{70}{cm} wavelength) penetrates canopy to trunks and large branches---directly relevant for buried wood structure detection. Data expected January 2026.

\subsection{Hyperspectral Imagery}

\begin{table}[h]
\centering
\caption{Hyperspectral Mission Comparison}
\label{tab:hyperspectral_comparison}
\begin{tabular}{p{2.5cm}ccccc}
\toprule
\textbf{Mission} & \textbf{Resolution} & \textbf{Bands} & \textbf{Wood Det.} & \textbf{Access} & \textbf{Status} \\
\midrule
EMIT & \SI{60}{m} & 285 & Moderate & Free & Active \\
PRISMA & \SI{30}{m} & 239 & Moderate & Free & Active \\
EnMAP & \SI{30}{m} & 224 & Moderate & Free & Active \\
\rowcolor{inprogress!30}
NASA SBG & \SI{30}{m} & $\sim$200+ & \textbf{High} & Free & 2027 \\
\bottomrule
\end{tabular}
\end{table}

\subsection{High-Resolution Optical}

\begin{table}[h]
\centering
\caption{Optical Imagery Comparison}
\label{tab:optical_comparison}
\begin{tabular}{p{2.5cm}ccccc}
\toprule
\textbf{Product} & \textbf{Resolution} & \textbf{Bands} & \textbf{Revisit} & \textbf{Access} & \textbf{Priority} \\
\midrule
Sentinel-2 & \SI{10-20}{m} & 13 & 5 days & Free & Medium \\
Landsat 8/9 & \SI{30}{m} & 11 & 16 days & Free & Low \\
\rowcolor{complete!30}
Planet NICFI & \SI{4.77}{m} & 4 & Monthly & \textbf{Free} & \textbf{HIGH} \\
Planet SuperDove & \SI{3}{m} & 8 & Daily & Commercial & Low \\
\bottomrule
\end{tabular}
\end{table}

\textbf{Key Finding:} Planet NICFI provides \textbf{free} \SI{4.77}{m} imagery for all tropical forests---6$\times$ improvement over Sentinel-2 for detailed scarp mapping.

%==============================================================================
\section{Prediction Models by Platform}
\label{sec:models}
%==============================================================================

\subsection{Model 1: Terrain-Based Scarp Probability (DEM)}

\subsubsection{Input Variables}
\begin{itemize}[noitemsep]
    \item Slope ($\theta$), normalized to 0--1
    \item TPI gradient ($\nabla TPI$), normalized
    \item Edge magnitude ($E_{mag}$), normalized
    \item Multi-scale TPI variance ($\sigma^2_{TPI}$), normalized
    \item Terrain roughness ($TRI$), normalized
\end{itemize}

\subsubsection{Model Formulation}

\begin{equation}
\label{eq:scarp_prob}
P_{scarp} = w_1 \cdot \theta_{norm} + w_2 \cdot \nabla TPI_{norm} + w_3 \cdot E_{mag} + w_4 \cdot \sigma^2_{TPI,norm} + w_5 \cdot TRI_{norm}
\end{equation}

\textbf{Default weights:} $w_1 = 0.30$, $w_2 = 0.25$, $w_3 = 0.20$, $w_4 = 0.15$, $w_5 = 0.10$

\subsubsection{Alternative DEM Products}

\begin{table}[h]
\centering
\caption{Expected Model Performance by DEM Product}
\label{tab:dem_model_performance}
\begin{tabular}{lccc}
\toprule
\textbf{DEM Product} & \textbf{Expected Precision} & \textbf{Expected Recall} & \textbf{Notes} \\
\midrule
NASADEM & 0.65--0.75 & 0.70--0.80 & Baseline \\
Copernicus & 0.65--0.75 & 0.70--0.80 & Similar to NASADEM \\
\textbf{FABDEM} & \textbf{0.75--0.85} & \textbf{0.80--0.90} & Canopy removal improves \\
WorldDEM DTM & 0.80--0.90 & 0.85--0.95 & Best but commercial \\
\bottomrule
\end{tabular}
\end{table}

\subsection{Model 2: Lidar-Based Height Validation (GEDI/ICESat-2)}

\subsubsection{Input Variables}
\begin{itemize}[noitemsep]
    \item Ground elevation from GEDI L2A ($z_{ground}$)
    \item Canopy height ($h_{canopy}$)
    \item Quality flag filtering (quality\_flag = 1, degrade\_flag = 0)
\end{itemize}

\subsubsection{Model Formulation}

Scarp height validation:
\begin{equation}
H_{scarp,lidar} = z_{terrace} - z_{floodplain}
\end{equation}

Cross-validation with DEM:
\begin{equation}
\Delta H = H_{scarp,DEM} - H_{scarp,lidar}
\end{equation}

\subsubsection{Platform Comparison}

\begin{table}[h]
\centering
\caption{Lidar Platform Performance for Height Estimation}
\label{tab:lidar_model_performance}
\begin{tabular}{lcccc}
\toprule
\textbf{Platform} & \textbf{Height RMSE} & \textbf{Tropical Perf.} & \textbf{Coverage} & \textbf{Recommendation} \\
\midrule
GEDI L2A & $\pm$\SI{1-2}{m} & Good & Sampling & \textbf{Primary} \\
ICESat-2 ATL08 & $\pm$\SI{2-4}{m} & Poor (NRMSE 67\%) & Sampling & Validation only \\
GEDI + ICESat-2 & $\pm$\SI{1.5}{m} & Moderate & Better & Combined approach \\
\bottomrule
\end{tabular}
\end{table}

\subsection{Model 3: Spectral Wood Detection (EMIT/Sentinel-2)}

\subsubsection{Input Variables}
\begin{itemize}[noitemsep]
    \item Cellulose absorption depth at \SI{2100}{nm}
    \item Lignin absorption depth at \SI{2270}{nm}
    \item NDVI anomaly (deviation from surrounding vegetation)
    \item Bare soil index (BSI)
\end{itemize}

\subsubsection{Model Formulation}

Cellulose Detection Index:
\begin{equation}
CDI = 1 - \frac{R_{2100}}{R_{continuum}}
\end{equation}

Combined spectral score:
\begin{equation}
S_{wood} = w_1 \cdot CDI + w_2 \cdot LDI + w_3 \cdot (1 - NDVI) + w_4 \cdot BSI
\end{equation}

\subsubsection{Platform Comparison}

\begin{table}[h]
\centering
\caption{Spectral Platform Performance for Wood Detection}
\label{tab:spectral_model_performance}
\begin{tabular}{lcccl}
\toprule
\textbf{Platform} & \textbf{Resolution} & \textbf{Wood Detection} & \textbf{Limitation} & \textbf{Use Case} \\
\midrule
EMIT & \SI{60}{m} & Moderate & Mixed pixels & Regional screening \\
PRISMA/EnMAP & \SI{30}{m} & Moderate & Limited coverage & Targeted analysis \\
Sentinel-2 & \SI{10-20}{m} & Low & No SWIR detail & Vegetation gaps only \\
\textbf{SBG (2027)} & \SI{30}{m} & High & Not yet available & Future primary \\
\bottomrule
\end{tabular}
\end{table}

\subsection{Model 4: SAR Structure Detection (BIOMASS/NISAR)}

\subsubsection{Input Variables}
\begin{itemize}[noitemsep]
    \item P-band backscatter (ESA BIOMASS)
    \item L-band backscatter (NISAR, ALOS-2)
    \item Polarimetric decomposition parameters
    \item InSAR coherence
\end{itemize}

\subsubsection{Model Formulation}

Biomass estimation (NISAR):
\begin{equation}
AGB = a \cdot \sigma^0_{HV,L} + b \cdot \sigma^0_{HV,S} + c
\end{equation}

Expected performance: $R^2 > 0.75$, RMSE \SIrange{15}{22}{Mg~ha^{-1}}

\subsubsection{Platform Comparison}

\begin{table}[h]
\centering
\caption{SAR Platform Performance for Structure Detection}
\label{tab:sar_model_performance}
\begin{tabular}{lcccc}
\toprule
\textbf{Platform} & \textbf{Wavelength} & \textbf{Penetration} & \textbf{Biomass Sat.} & \textbf{Status} \\
\midrule
Sentinel-1 (C-band) & \SI{5.6}{cm} & Canopy top & \SI{50}{Mg~ha^{-1}} & Active \\
ALOS-2 (L-band) & \SI{24}{cm} & Mid-canopy & \SI{100}{Mg~ha^{-1}} & Active \\
NISAR (L+S) & \SI{24}{cm} + \SI{10}{cm} & Mid + upper & \SI{150}{Mg~ha^{-1}} & July 2025 \\
\textbf{BIOMASS (P-band)} & \SI{70}{cm} & \textbf{Trunks} & $>$\SI{300}{Mg~ha^{-1}} & \textbf{Data 2026} \\
\bottomrule
\end{tabular}
\end{table}

\subsection{Model 5: Integrated Multi-Platform Prediction}

\subsubsection{Ensemble Approach}

\begin{equation}
P_{final} = \alpha \cdot P_{terrain} + \beta \cdot P_{lidar} + \gamma \cdot P_{spectral} + \delta \cdot P_{SAR}
\end{equation}

Where weights $\alpha, \beta, \gamma, \delta$ are determined by:
\begin{enumerate}[noitemsep]
    \item Individual platform accuracy (from validation)
    \item Data availability at location
    \item Uncertainty of each prediction
\end{enumerate}

\subsubsection{Decision Framework}

\begin{table}[h]
\centering
\caption{Multi-Platform Decision Matrix}
\label{tab:decision_matrix}
\begin{tabular}{lccccc}
\toprule
\textbf{Scenario} & \textbf{DEM} & \textbf{GEDI} & \textbf{EMIT} & \textbf{SAR} & \textbf{Confidence} \\
\midrule
All agree high & High & High & High & High & Very High \\
Terrain only & High & Gap & Low & Low & Moderate \\
Terrain + lidar & High & High & Gap & Low & High \\
Spectral anomaly & Low & Gap & High & Mod & Moderate \\
SAR structure & Mod & Gap & Gap & High & Moderate \\
\bottomrule
\end{tabular}
\end{table}

%==============================================================================
\section{Science Traceability Matrix}
\label{sec:stm}
%==============================================================================

\begin{landscape}
\begin{table}[p]
\centering
\caption{Full Science Traceability Matrix}
\label{tab:stm_full}
\small
\begin{tabular}{p{1.5cm}p{4cm}p{2.5cm}p{2.5cm}p{3cm}p{2cm}c}
\toprule
\textbf{ID} & \textbf{Science Question} & \textbf{Measurement} & \textbf{Data Product} & \textbf{Variable} & \textbf{Validation} & \textbf{Feasibility} \\
\midrule
\multicolumn{7}{l}{\textbf{Scarp Detection}} \\
STM-1.1 & Where are terrace scarps? & Location & NASADEM, FABDEM & Slope $>$\ang{15} & GPS survey & \HIGH \\
STM-1.2 & Where are terrace scarps? & Location & DEM & TPI gradient & GPS survey & \HIGH \\
STM-1.3 & Where are terrace scarps? & Location & DEM & Edge magnitude & Field mapping & \HIGH \\
STM-1.4 & What is scarp probability? & Probability (0--1) & Derived & Composite score & Confusion matrix & \HIGH \\
\midrule
\multicolumn{7}{l}{\textbf{Geometry Measurement}} \\
STM-2.1 & What is scarp height? & Height (m) & GEDI L2A, DEM & Ground elevation & Rangefinder & \HIGH \\
STM-2.2 & What is scarp length? & Length (m) & Scarp vectors & Centerline & Field measure & \HIGH \\
STM-2.3 & What is scarp orientation? & Azimuth (\ang{}) & Scarp vectors & Direction & Compass & \HIGH \\
\midrule
\multicolumn{7}{l}{\textbf{Wood Detection}} \\
STM-3.1 & Is wood spectrally detectable? & Reflectance & EMIT L2A & Cellulose \SI{2100}{nm} & Spectrometer & \MODERATE \\
STM-3.2 & Is organic matter elevated? & Index & EMIT L2A & Lignin \SI{2270}{nm} & Lab chemistry & \MODERATE \\
STM-3.3 & Are vegetation gaps present? & NDVI & Sentinel-2, NICFI & NDVI anomaly & Visual & \MODERATE \\
STM-3.4 & Is wood structure detectable? & Backscatter & ESA BIOMASS & P-band $\sigma^0$ & TBD & \MODERATE \\
\midrule
\multicolumn{7}{l}{\textbf{Carbon Quantification}} \\
STM-4.1 & What is wood volume? & \unit{m^3~m^{-1}} & Field data & Cross-section & Excavation & \LOW \\
STM-4.2 & What is wood density? & \unit{kg~m^{-3}} & Lab analysis & Dry weight/vol & Standard & \LOW \\
STM-4.3 & What is carbon fraction? & \% & Lab analysis & Elemental C & CHN analyzer & \LOW \\
STM-4.4 & What is preservation state? & \% & Lab analysis & Mass loss & Reference & \LOW \\
STM-4.5 & What is burial age? & ka BP & Lab analysis & $^{14}$C / OSL & AMS dating & \LOW \\
\bottomrule
\end{tabular}
\end{table}
\end{landscape}

%==============================================================================
\section{Data Product Status and Tracking}
\label{sec:tracking}
%==============================================================================

\subsection{Current Acquisition Status}

\begin{table}[h]
\centering
\caption{Data Product Acquisition Status}
\label{tab:status}
\begin{tabular}{llccc}
\toprule
\textbf{Product} & \textbf{Source} & \textbf{Status} & \textbf{Priority} & \textbf{Action} \\
\midrule
\multicolumn{5}{l}{\textit{Complete}} \\
NASADEM & NASA EarthData & \status{complete} Complete & -- & -- \\
Copernicus DEM & AWS/ESA & \status{complete} Complete & -- & -- \\
Terrain derivatives & Derived & \status{complete} Complete & -- & -- \\
\midrule
\multicolumn{5}{l}{\textit{In Progress}} \\
GEDI L4A & NASA LP DAAC & \status{inprogress} Downloading & HIGH & Monitor \\
EMIT L2A & NASA LP DAAC & \status{inprogress} Downloading & HIGH & Monitor \\
\midrule
\multicolumn{5}{l}{\textit{Priority Pending}} \\
\textbf{FABDEM} & OpenTopography & \status{pending} Pending & \textbf{HIGH} & Download now \\
\textbf{Planet NICFI} & Planet & \status{pending} Pending & \textbf{HIGH} & Register \\
ALOS-2 PALSAR & JAXA/GEE & \status{pending} Pending & Medium & Download \\
\midrule
\multicolumn{5}{l}{\textit{Future Missions}} \\
NISAR & NASA/ISRO & \status{pending} July 2025 & HIGH & Monitor \\
ESA BIOMASS data & ESA & \status{pending} Jan 2026 & HIGH & Monitor \\
NASA SBG & NASA & \status{pending} 2027 & Medium & Monitor \\
\bottomrule
\end{tabular}
\end{table}

%==============================================================================
\section{Validation Plan}
\label{sec:validation}
%==============================================================================

\subsection{Validation Hierarchy}

\begin{enumerate}
    \item \textbf{V1: Internal consistency} -- DEM product inter-comparison
    \item \textbf{V2: Geometric accuracy} -- GEDI + GPS for scarp location/height
    \item \textbf{V3: Thematic accuracy} -- Field verification of scarp detection
    \item \textbf{V4: Biophysical validation} -- Wood presence at detected scarps
    \item \textbf{V5: Quantitative validation} -- Carbon estimates vs. lab analysis
\end{enumerate}

\subsection{Performance Metrics}

\begin{table}[h]
\centering
\caption{Validation Metrics and Targets}
\label{tab:validation_metrics}
\begin{tabular}{llcc}
\toprule
\textbf{Metric} & \textbf{Formula} & \textbf{Target} & \textbf{Current} \\
\midrule
Detection accuracy & $(TP + TN) / (TP + TN + FP + FN)$ & $>$80\% & TBD \\
Precision & $TP / (TP + FP)$ & $>$70\% & TBD \\
Recall & $TP / (TP + FN)$ & $>$80\% & TBD \\
Height RMSE & $\sqrt{\sum(h_{pred} - h_{obs})^2 / n}$ & $<$\SI{3}{m} & TBD \\
Carbon error & $|C_{pred} - C_{obs}| / C_{obs}$ & $<$50\% & TBD \\
\bottomrule
\end{tabular}
\end{table}

%==============================================================================
\section{Timeline and Milestones}
\label{sec:timeline}
%==============================================================================

\begin{table}[h]
\centering
\caption{Project Milestones}
\label{tab:milestones}
\begin{tabular}{lll}
\toprule
\textbf{Phase} & \textbf{Activities} & \textbf{Deliverables} \\
\midrule
\textbf{Phase 1: Data Acquisition} & & \\
\quad Current & GEDI, EMIT downloading & Raw data archive \\
\quad Immediate & FABDEM, Planet NICFI & Enhanced DEM, high-res optical \\
\quad Near-term & Sentinel-2, ALOS-2 & Spectral indices, SAR data \\
\midrule
\textbf{Phase 2: Analysis} & & \\
\quad Processing & Multi-platform terrain analysis & Scarp probability maps \\
\quad Integration & Cross-platform validation & Ensemble predictions \\
\quad Comparison & Platform performance metrics & Model evaluation report \\
\midrule
\textbf{Phase 3: Field Campaign} & & \\
\quad Verification & Stratified site visits (n=30+) & Validation dataset \\
\quad Sampling & Wood/sediment collection & Lab samples \\
\midrule
\textbf{Phase 4: Publication} & & \\
\quad Analysis & Carbon estimation, uncertainty & Final estimates \\
\quad Writing & Manuscript preparation & Submitted paper \\
\bottomrule
\end{tabular}
\end{table}

%==============================================================================
\section{Expert Review Notes}
\label{sec:review}
%==============================================================================

\textit{[This section reserved for expert comments and revisions]}

\vspace{2cm}

\noindent\textbf{Reviewer Comments:}

\vspace{5cm}

\noindent\textbf{Action Items:}

\vspace{3cm}

%==============================================================================
% REFERENCES
%==============================================================================

\section*{References}
\addcontentsline{toc}{section}{References}

\begin{enumerate}[label={[\arabic*]}, leftmargin=*]
    \item National Academies of Sciences, Engineering, and Medicine. (2019). \textit{Negative Emissions Technologies and Reliable Sequestration: A Research Agenda}. Washington, DC: The National Academies Press. \url{https://doi.org/10.17226/25259}

    \item National Academies of Sciences, Engineering, and Medicine. (2018). \textit{Thriving on Our Changing Planet: A Decadal Strategy for Earth Observation from Space}. Washington, DC: The National Academies Press.

    \item Zeng, N. (2008). Carbon sequestration via wood burial. \textit{Carbon Balance and Management}, 3:1. \url{https://doi.org/10.1186/1750-0680-3-1}

    \item Hawker, L., et al. (2022). A 30 m global map of elevation with forests and buildings removed. \textit{Environmental Research Letters}, 17(2), 024016.

    \item Dubayah, R., et al. (2020). The Global Ecosystem Dynamics Investigation: High-resolution laser ranging of the Earth's forests and topography. \textit{Science of Remote Sensing}, 1, 100002.

    \item Le Toan, T., et al. (2011). The BIOMASS mission: Mapping global forest biomass to better understand the terrestrial carbon cycle. \textit{Remote Sensing of Environment}, 115(11), 2850-2860.

    \item Rosen, P.A., et al. (2017). The NASA-ISRO SAR (NISAR) mission dual-band radar instrument preliminary design. \textit{IEEE International Geoscience and Remote Sensing Symposium}, 3832-3835.
\end{enumerate}

\end{document}
